\chapter{绪论}\label{ch:intro}

\section{研究背景与意义}\label{sec:background}

近年来，随着第五代移动通信技术（The Fifth Generation Mobile Communication Technology, 5G）在全球范围内的规模化商用部署，移动通信技术在提升信息传输效率、拓展应用场景以及推动社会信息化进程等方面发挥了重要作用\cite{YouXiaoHu20145gyidongtongxinfazhanqushiyuruoganguanjianjishu, ahmed2018survey}。依托大规模多输入多输出（Massive Multiple-Input and Multiple-Output, mMIMO）、毫米波（Millimeter Wave, mmWave）通信以及超密集组网（Ultra-Dense Network, UDN）等关键技术，5G网络在峰值速率、频谱效率和连接密度等性能指标上相较第四代移动通信技术（The Fourth Generation Mobile Communication Technology, 4G）实现了数量级的提升，初步支撑了增强型移动宽带（Enhanced Mobile Broadband, eMBB）、大规模机器通信（Massive Machine-Type Communications, mMTC）以及超可靠低时延通信（Ultra-Reliable and Low-Latency Communications, URLLC）这三大应用场景的需求\cite{marzetta2016fundamentals,marzetta2015massive,larsson2014massive,bennis2018ultrareliable}。
然而，远程医疗、无人驾驶和沉浸式虚拟现实等技术的蓬勃发展，对无线链路的速率、时延和能效等要求越来越高，使得现有技术的局限性逐渐显现出来，人们开始了对第六代移动通信技术（The Sixth Generation Mobile Communication Technology, 6G）的预研\cite{zhao20196g,ZhangPing20196gyidongtongxinjishuzhanwang,you20216g}。

\begin{figure}[!htbp]
	\centering
	% 第一张子图
	\begin{subfigure}[b]{0.50\linewidth} % 设置子图宽度，[b]表示底部对齐
		\centering
		\includegraphics[width=\linewidth]{intro/intro_6G_usage_scenarios.pdf}
		\caption{6G典型应用场景}
		\label{fig:intro_6G_usage_scenarios}
	\end{subfigure}% % 这里的 % 很有必要，防止产生多余空格导致换行
	\hfill % 在两个子图之间自动填充间距
	% 第二张子图
	\begin{subfigure}[b]{0.46\linewidth}
		\centering
		\includegraphics[width=\linewidth]{intro/intro_6G_performance.pdf}
		\caption{6G与5G性能指标对比}
		\label{fig:intro_6G_performance}
	\end{subfigure}
	\vspace{-0.5cm}
	\caption{IMT-2030中提出的6G愿景与性能指标\cite{ITU2023M2160}}
	\label{fig:intro_6G_imt}
\end{figure}

2023年11月，国际电信联盟（International Telecommunication Union, ITU）发布了6G建议书\cite{ITU2023M2160}，书中明确提出6G网络的六大典型应用场景，如图~\ref{fig:intro_6G_imt}(a)所示。其中包含对5G的原有场景提出了更高的要求，分别是：沉浸式通信、大规模通信和极致可靠和低延时通信，另外引入了通信感知一体化（Integrated Sensing and Communication, ISAC）、通信与人工智能（Artificial Intelligence, AI）融合和泛在连接三大全新应用范式。为了实现这些愿景，ITU为6G网络提出了15个具体的性能指标维度要求，如图~\ref{fig:intro_6G_imt}(b)所示。具体而言，九大增强型能力指标在5G原有能力上拓展，数据速率、频谱效率和能量效率均要求数倍于5G性能；六大全新能力涵盖全域覆盖、感知能力、定位能力、AI能力、可持续性和互操作性，要求6G网络向智能化、一体化发展。

然而，5G系统中广泛使用的技术已经不能很好地适应这些新需求。一方面，传统的mMIMO技术由有源相控阵实现，通过在基站侧和用户侧配置多根天线，可以带来信道容量的提升，且增益正比于天线数量\cite{asplund2020advanced,bjornson2016massive,yin2020addressing}。为实现6G中更高的频谱效率要求，所需的天线数量将剧烈增长，这不仅受到系统物理尺寸和复杂度的制约，大量的射频链路还会引起不可忽视的成本和功耗问题\cite{dohler2010cooperative}。另一方面，现有的蜂窝网络普遍受限于既定的物理传播环境，依靠更高的建站密度提升非视距（Non-Line-of-Sight, NLoS）覆盖，实现“全域覆盖”将进一步推高网络的建设和运营成本。在6G对高性能与低功耗双重要求的背景下，需要打破传统的通信架构，由“适应信道”的旧范式转变为“主动改变信道”的新设计范式，力求实现高效节能通信。

因此，如何构建智能电磁环境成为6G技术演进的重要突破方向。可重构天线系统（Reconfigurable Antenna System, RAS）这一新型体系结构，能够主动调控电磁波的传播行为，实现“重构信道”，已经成为6G物理层系统架构中的重要研究方向之一。广义而言，RAS是指一类能够在相位响应、幅度响应或空间分布等维度实现动态调控的天线阵列或电磁表面，其核心目标是在尽量降低硬件复杂度的前提下，实现对电磁波传播路径和特性的主动控制，有效解决传统mMIMO技术在成本、能效与部署方面面临的挑战。根据系统部署位置及重构方式的不同，现有RAS相关研究可大致划分为三类具有代表性的技术路径：第一类为部署于发射端，用以取代或简化有源相控阵的动态超表面天线（Dynamic Metasurface Antenna, DMA）\cite{smith2017analysis}；第二类为同样部署于发射端，但通过调整天线空间位置可实现大尺度信道重构的捏合天线系统（Pinching Antenna Systems, PASS）\cite{ding2025flexibleantenna}；第三类为部署在基站（Base Station, BS）与用户设备（User Equipment, UE）之间，用于构建视距链路（Line-of-Sight, LoS），改善覆盖条件的可重构智能表面（Reconfigurable Intelligent Surface, RIS）\cite{liu2021reconfigurable}。下文将对上述三类可重构天线技术的基本原理与典型特征进行简要介绍。

其中，动态超表面天线侧重于基站侧的体系结构革新，不同于传统相控阵高度依赖射频链路或移相器网络进行相位调节，DMA通过在波导或谐振腔表面集成大量包含正-本征-负（Positive-Intrinsic-Negative, PIN）二极管、变容二极管等可调谐元件的超材料单元，实现对电磁波幅度和相位的直接调控。借助对单元谐振状态的独立控制，DMA能够显著减少天线阵列中射频链路的数量，从而降低硬件复杂度和系统功耗，为超大规模天线阵列的工程实现提供了可行路径。然而，这种以波导馈电和单元辐射为特征的结构形式，也使DMA的信号传输过程不可避免地受到波导传播效应及单元耦合特性的共同影响，其可实现的波束增益及性能边界有别于传统相控阵。因此，有必要从理论层面对DMA的波束赋形机制及其性能极限进行系统分析。

在此基础上，捏合天线系统进一步从空间层面拓展了可重构自由度。捏合天线（Pinching Antenna, PA）又称夹子天线，介质波导上被夹子夹住的地方，介电常数会发生变化，电磁波便会辐射出来。利用这一性质，将DMA中波导上的超材料单元替换成这种可以沿着波导移动的“夹子”，使其具有位置可调的新自由度。该特性使捏合天线系统能够主动选择有利的空间位置，以规避深度衰落的区域并降低路径损耗，同时为高精度定位与感知提供更为丰富的空间信息。但实际部署中天线位置的移动往往伴随着馈电波导长度的改变，使信号在波导中传输和自由空间传播这两个过程存在紧密的耦合关系。上述因素使捏合天线的位置优化复杂化，需要针对性进行性能分析与算法设计。因此，有必要建立符合物理传播特性的PASS传输模型，定量刻画空间自由度引入所带来的性能增益，并进一步探索PA位置信息辅助下的通感一体化波束赋形机制。

与DMA和PASS侧重于发射机自身的重构不同，RIS将重构对象拓展至无线传播环境本身。RIS是一种由大量的亚波长人工结构单元组成的平面阵列，具备调控电磁波幅度、相位或极化的能力。与DMA类似，RIS单元通常由电阻、电容、电感和二极管等无源元件构成，调节加载在单元上的电压值，能产生不同的响应。于是，通过对各单元反射相位的协同调控，RIS能够重塑电磁波的传播路径，实现信号在空间中的定向反射与聚焦。由于RIS能够以极低的功耗人为构建绕过障碍物的视距链路，它很适合作为部署于基站与用户之间的辅助通信节点，在提升覆盖的同时降低系统能耗。然而，RIS作为独立于收发机之外的即插即用型系统，通常难以与基站或终端协同完成信道估计过程；同时，其自身缺乏独立的射频链路，难以直接感知获取信道状态信息（Channel State Information, CSI），使得依赖完备CSI的传统波束赋形与优化方法难以适用。因此，如何设计无需信道状态信息的低复杂度波束赋形算法，并在算法与工程实现之间建立可行的衔接路径，构建可实际部署的原型系统，是RIS技术走向商用的关键前提。

尽管DMA、PASS与RIS在具体实现形式和调控维度上各有侧重，但三者的核心目标一致，即通过引入新的信道调控自由度，实现对电磁传播环境的精细控制。无论是简化射频结构、重构天线空间位置，还是调控传播环境，其最终效果均体现为波束赋形能力的重塑及系统性能的提升。因此，围绕波束赋形机理、实际场景中的性能分析以及低复杂度算法实现开展研究，构成了可重构天线系统在6G物理层设计中的共同理论基础。下一节将分别综述三类可重构天线技术的代表性研究成果并总结目前面临的挑战。


\section{国内外研究现状与挑战}\label{sec:current_research_and_challenge}

\subsection{动态超表面天线研究现状与挑战}

如前文所述，动态超表面天线的核心是引入了包含可调谐元件的超材料单元。超材料这一概念在20世纪60年代由前苏联理论物理学家维克托·韦谢拉戈提出，他预言了磁导率和介电常数同为负值的左手材料的存在。他在研究中指出，当电磁波在这种介质中传播时，其电场矢量、磁场矢量与波矢量三者之间的方向关系，不再符合经典电磁学所描述的右手定则，而表现为左手关系\cite{veselago1968electrodynamics}。其传播相速度的方向与能量流动的方向恰好相反，将产生“负折射”的现象。这一现象为控制电磁波的传播开拓了新的可能性\cite{engheta2006metamaterials}。然而，受限于实验技术和制作工艺，这一构想被长期搁置，在当时并未引起足够的关注。

20世纪末，随着材料科学的不断进步，研究者们通过在材料上设计精细的微观结构，逐步制备出具备负磁导率和负介电常数的人工材料，并在此基础上开展对电磁波调控的研究。1996年，英国的J. B. Pendry教授提出了设计负介电常数结构的思路，利用周期性金属细线阵列结构，将金属的有效等离子体频率大幅降低至微波频段，实现了在该频段下具有等效负介电常数的人工超材料\cite{pendry1996extremely,pendry1998low}。随后，这一团队在1999年进一步提出了实现负磁导率超材料的结构方案：采用亚波长尺寸的金属开口谐振环组成周期性阵列，通过人为设计可使该结构表现出负磁导率特性\cite{pendry1999magnetism}。Pendry 教授主要贡献在于提出了较为完善的超材料设计理论，直至2001年，杜克大学的D. R. Smith教授团队首次采用印刷电路板（Printed Circuit Board, PCB）工艺成功制备出实物，并在实验中观测到了双负现象，验证了超材料相关理论的正确性\cite{smith2000composite,shelby2001experimental}。

相较于结构复杂、加工难度高的传统三维体态超材料，二维超表面以其低剖面、易集成的特点，成为一种新兴的超材料形态，能够有效简化制造工艺并降低集成复杂度\cite{chen2016review}。







\noindent\rule{\linewidth}{2pt}

随着制造工艺与电磁理论的不断演进，三维的体态超材料逐渐向低剖面、易集成的二维超表面形态过渡。正是在这种通过人工微结构重塑电磁波行为的长期积淀下，D. R. Smith 团队顺理成章地将超材料技术引入了微波天线的设计中，开创性地提出了一种颠覆传统射频架构的波导馈电动态超表面天线（Dynamic Metasurface Antenna, DMA）模型\cite{smith2017analysis}。



这种全新的架构将微带波导或二维谐振腔作为信号传输的骨架，在其表面密集嵌刻了大量亚波长尺寸的辐射缝隙。更为精妙的是，每个辐射单元内部均集成了如PIN二极管等有源可调器件。在发射状态下，基带信号仅需由极少量的射频链路注入波导。当导波在结构内部向前传播时，受数字指令控制的二极管会动态改变各个辐射单元的电磁谐振状态，使导波以漏波（Leaky-wave）的形式耦合至自由空间。这种设计直接在天线的物理辐射端完成了波束赋形的幅度与相位重构，彻底剥离了传统相控阵对庞大移相器网络和密集射频链路的依赖。凭借其低剖面、低功耗以及极简馈电的颠覆性优势，DMA迅速成为6G超大规模天线系统演进路径上极具潜力的研究焦点。



在向6G演进的进程中，天线阵列规模的持续扩张与硬件能耗之间的矛盾已成为无线通信物理层设计中难以回避的核心痛点。现有的全数字或混合预编码架构，本质上依赖于庞大的射频链路群与精密（且高功耗）的移相器网络来实现大规模MIMO的波束赋形。然而，当天线数量达到数百乃至上千级别时，这种“堆叠硬件”的传统相控阵设计思路在成本控制、馈电网络布线以及散热管理上均面临着难以逾越的物理瓶颈\cite{bjornson2017massive}。要从根本上打破这一僵局，就必须在天线的底层物理机制上寻求突破，而近年来超材料（Metamaterials）在电磁波调控领域的快速发展，恰好为重构天线阵列形态提供了全新的技术土壤\cite{cui2010metamaterials}。
通过将超材料的亚波长周期性结构与微波传输线相结合，杜克大学D. R. Smith团队首次提出了一种全新的天线架构——动态超表面天线（Dynamic Metasurface Antenna, DMA）\cite{smith2017analysis}。这种架构从根本上颠覆了传统天线的馈电与辐射机制。在具体实现上，DMA通常由一维微带波导或二维谐振腔构成主体，并在其表面密集排列大量亚波长尺寸的辐射单元。这里的关键在于，每个辐射单元不仅是物理意义上的“天线缝隙”，更是一个加载了PIN二极管或变容二极管的可调谐电磁谐振器\cite{shlezinger2019dynamic}。
从信号传输的视角来看，DMA的工作机理非常独特。在发射端，携带基带信息的电磁波只需由极少量的射频链路注入到微带波导的起始位置。当导波在波导内部向前传输时，它会依次经过表面的各个辐射单元，并通过电磁耦合以漏波（Leaky-wave）的形式向外辐射。由于表面每个单元的谐振状态均受其内部二极管偏置电压的独立控制，基站只需通过极低功耗的数字指令改变这些电压状态，就能直接在物理层面上对辐射出去的电磁波幅度与相位进行联合重构。这种将复杂的基带波束赋形操作下沉至天线物理辐射层的设计，不仅使系统摆脱了对密集射频链路的依赖，还在剖面尺寸和功耗上展现出传统相控阵难以企及的优势，为面向未来的超大规模智能天线设计开辟了一条极具工程潜力的技术路线。

如前文所述，为了满足6G网络对极高频谱效率与全域覆盖的严苛要求，基站端配置的天线阵列规模正向着超大阵列（Ultra-Massive MIMO）加速演进。然而，传统的全数字或混合相控阵天线架构通常需要为大量的天线振子配备独立的射频链路（RF Chains）与高精度的移相器网络。随着阵列规模的成倍扩张，庞大的射频硬件不仅导致系统成本与功耗呈指数级激增，其复杂的微带馈电布线也在高频段（如毫米波系统）面临着严重的物理空间与电磁互耦瓶颈\cite{bjornson2017massive}。为了打破这一“功耗-成本-尺寸”的工程瓶颈，研究人员开始寻求天线底层物理架构的根本性变革，而超材料（Metamaterials）技术的突破恰好为此提供了契机。

超材料是由亚波长尺度的人工微结构按特定空间排列方式构成的新型电磁材料。它突破了自然材料的电磁属性限制，能够实现负折射率、完美吸收以及异常反射等奇异的电磁调控现象\cite{cui2010metamaterials}。随着该技术的不断成熟，杜克大学（Duke University）的 D. R. Smith 团队将其与传统导波结构相结合，于近年来开创性地提出了一种基于波导馈电的动态超表面天线（Dynamic Metasurface Antenna, DMA）架构\cite{smith2017analysis}。

与传统相控阵复杂的树状馈电网络或自由空间馈电不同，DMA 巧妙地在一维微带波导或二维谐振腔体的表面，密集刻蚀了大量亚波长尺寸的辐射单元（Metamaterial Elements）以形成物理阵列。每一个辐射单元本质上是一个可调谐的电磁谐振器，其内部集成了正-本征-负（PIN）二极管或变容二极管（Varactor Diodes）等有源固态半导体器件\cite{shlezinger2019dynamic}。在DMA的工作过程中，基带信号仅需通过极少数的射频链路馈入波导的起始端。当电磁波在波导内部向前传播时，会与表面的超材料辐射单元发生电磁耦合，从而形成向自由空间辐射的漏波（Leaky-wave）。基站控制器通过外部指令动态调节各单元内二极管的偏置电压，即可独立改变每个超材料单元的谐振频率，进而在天线的物理辐射层面实现对电磁波幅度与相位的直接重构。这种创新架构将波束赋形操作从高功耗的数字基带转移到了天线物理层，极大地压缩了所需的射频链路数量，凭借其低剖面、低功耗的显著优势，成为了学术界公认的下一代大规模天线演进方向之一。



超材料（Metamaterials）是由亚波长尺度的人工微结构按特定空间排列方式构成的新型电磁材料。它突破了自然材料的电磁属性限制，能够实现负折射率、完美吸收以及异常反射等奇异的电磁现象\cite{cui2010metamaterials}。随着超材料技术的不断成熟，其在天线设计领域的应用逐渐成为学术界与工业界的研究热点。传统的大规模MIMO和毫米波通信系统高度依赖于全数字或混合相控阵架构，这类架构通常需要为每根天线或子阵列配备独立的射频链路（RF Chains）与高精度移相器。随着天线规模的急剧增加，系统的硬件成本、馈电网络布线难度以及射频功耗均呈现出难以接受的线性甚至指数级增长\cite{bjornson2017massive}，成为了制约6G大阵列系统走向实用化的核心瓶颈。

\noindent\rule{\linewidth}{2pt}

超材料（Metamaterial）是一类通过人工设计的周期性或非周期性微结构单元构成的复合材料，其电磁特性主要取决于单元结构的几何形状与空间排列方式，而非材料的本征属性。这一独特特性使得超材料能够展现出天然材料所不具备的电磁响应行为，例如负折射率、完美吸收、电磁波极化调控等奇异现象。自21世纪初Pendry等人关于完美透镜的开创性工作以来，超材料在电磁波调控领域的研究取得了蓬勃发展，并逐渐从基础物理研究走向工程应用。在天线设计领域，超材料的应用为突破传统天线性能瓶颈提供了全新思路。传统天线设计往往受限于材料的固有属性和几何尺寸，而超材料天线通过精心设计单元结构，可以实现对电磁波传播方向、极化状态、频率响应等特性的灵活调控，从而显著提升天线的增益、带宽、扫描范围等关键性能指标。

在超材料天线的发展历程中，动态超表面天线（Dynamic Metasurface Antenna, DMA）的提出标志着一个重要里程碑。2017年，Smith团队在文献中首次系统地提出了波导馈电DMA的概念与理论模型。与传统相控阵天线依赖大量射频链路和移相器实现波束扫描不同，DMA采用了一种创新性的架构设计：将亚波长尺度的超材料辐射单元嵌入波导或腔体结构中，通过外部控制信号调节每个单元的电磁响应特性，从而实现对辐射波前的程序化调控。在Smith团队提出的波导馈电DMA模型中，每条波导上均匀分布着若干超材料辐射单元，输入信号经由波导传输并逐级馈电至各个单元，最终由单元向自由空间辐射。这一架构的核心优势在于，它仅需要与波导数量相等的射频链路，而非与辐射单元数量相等的射频链路，从而大幅降低了系统的硬件复杂度和功耗成本。

从物理建模角度，Smith团队的研究揭示了DMA辐射单元的频率响应呈现出典型的洛伦兹谐振特性。具体而言，每个超材料单元可视为一个可调谐的谐振电路，其频率响应可通过调节单元的谐振频率进行外部控制。对于窄带通信系统，单元的响应可以简化为洛伦兹约束形式$q = (j + e^{j\phi})/2$，其中$\phi$为可配置的相位偏移参数。这一数学形式揭示了DMA设计中面临的一个关键挑战：单元的幅度响应与相位响应存在固有的耦合关系，而非传统相控阵中相互独立的幅度和相位控制。这种耦合特性使得DMA的波束赋形优化问题具有独特的数学结构，既不同于传统的数字预编码问题，也不同于基于单位模约束的移相器优化问题。Smith团队的工作不仅建立了DMA的理论分析框架，还深入探讨了互耦效应、波导衰减、单元间距等实际因素对系统性能的影响，为后续DMA在无线通信领域的应用研究奠定了坚实基础。



\noindent\rule{\linewidth}{2pt}



为了突破传统相控阵天线的物理与功耗瓶颈，杜克大学（Duke University）的 D. R. Smith 团队于近年来开创性地提出了一种基于波导馈电的动态超表面天线（Dynamic Metasurface Antenna, DMA）架构\cite{smith2017analysis}。与传统相控阵复杂的微带馈电网络或自由空间馈电不同，DMA 在一维微带波导或二维谐振腔体的表面密集刻蚀大量亚波长尺寸的辐射单元（Metamaterial Elements）以形成物理阵列。每一个辐射单元本质上是一个可调谐的电磁谐振器，其内部集成了正-本征-负（PIN）二极管或变容二极管（Varactor Diodes）等有源固态半导体器件\cite{shlezinger2019dynamic}。

在DMA的工作过程中，基带信号仅需通过极少数的射频链路馈入波导的起始端。当电磁波在波导内部向前传播时，会与表面的超材料辐射单元发生电磁耦合，从而形成向自由空间辐射的漏波（Leaky-wave）。基站控制器通过外部的现场可编程逻辑门阵列（FPGA）或微控制器发出数字指令，动态调节各单元内二极管的偏置电压，即可独立改变每个超材料单元的谐振频率，进而实现对辐射电磁波幅度与相位的直接重构。这种创新的物理架构巧妙地将波束赋形操作从高功耗的基带数字预编码或射频移相器网络，转移到了天线阵列的物理辐射层面。它不仅极大地压缩了多天线系统所需的射频链路数量，还凭借其低剖面、低成本、低功耗的显著优势，为6G超大规模天线阵列的工程部署提供了一种极具潜力的革命性替代方案。

\noindent\rule{\linewidth}{2pt}

\begin{figure}[!htbp]
	\centering
	\includegraphics[width=0.5\linewidth]{intro/intro_DMA_waveguide_loss.pdf}
	\caption{动态超表面天线结构及工作原理示意图\cite{yoo2019enhancing}}
	\label{fig:intro_DMA_waveguide_loss}
\end{figure}



SIW

介绍第一次提出


介绍洛伦兹约束

从极化率的角度详细分析为什么会相依能力不足


介绍波束赋形算法

介绍性能分析现状

介绍挑战


这一部分需要展示对DMA底层机理的深刻理解，并精准地指出当前理论研究与工程实现之间的脱节，从而为你后续提出“离散算法”和“量化分析”做铺垫。

以下是为你撰写的 **1.2.1 动态超表面天线 (DMA) 研究现状**，语言风格保持学术严谨，你可以根据实际引用的文献编号进行调整。

---

 1.2.1 动态超表面天线 (DMA) 研究现状

作为一种极具潜力的新型收发机架构，动态超表面天线（DMA）通过在波导或空腔结构上集成大量可调超材料单元，实现了对电磁波辐射特性的动态调控。与传统相控阵天线依赖移相器和独立射频链路不同，DMA利用超表面的物理特性进行信号处理，极大地简化了硬件架构。

**（1）DMA架构与洛伦兹谐振模型**
DMA的基本概念最早由杜克大学的D. R. Smith教授团队系统性提出[1]。其典型架构通常由微带线或波导作为馈电结构，表面周期性或准周期性地排列着亚波长超材料谐振单元。当导波在馈电结构中传输时，通过外部控制信号（如直流偏置电压）改变加载在单元上的电子元件（如变容二极管或PIN二极管）的状态，从而调节每个单元的辐射振幅和相位，最终在自由空间合成所需的波束方向图。

在理论建模方面，Smith等人提出了基于**洛伦兹（Lorentzian）谐振**的极化率模型来描述DMA单元的电磁响应[2]。在该模型中，第  个超材料单元的磁极化率  通常表示为：



其中， 为工作频率， 为单元的谐振频率， 为振荡强度， 为阻尼系数。通过调节电路参数改变谐振频率  或阻尼 ，即可改变单元对导波的耦合强度和辐射相位。这一模型揭示了DMA的一个关键物理特性：**幅度与相位是强耦合的**，即在改变相位的过程中，单元的辐射幅度也会随之遵循洛伦兹曲线发生非线性变化，这与传统相控阵中幅度相位独立调控的机制截然不同。

**（2）基于连续相位假设的波束赋形算法研究**
基于上述架构，早期的DMA波束赋形研究主要聚焦于全息波束赋形理论。其核心思想是将所需的波束方向图视为“目标物波”，将波导内的传输场视为“参考波”，通过优化超表面单元的阻抗分布来重建目标波前。为了简化优化问题，现有文献多采用**松弛（Relaxation）策略**，即假设超材料单元的谐振频率或相位响应可以在一定范围内**连续变化**。
在此假设下，研究者们将DMA的波束赋形建模为连续域上的数学优化问题，并利用最小二乘法（Least Squares）、梯度下降法或基于格林函数的迭代算法求解最优的单元配置权重[3-4]。这些基于连续性假设的算法在理论仿真中展示了DMA在多波束生成、干扰抑制及高维MIMO传输方面的巨大潜力。

**（3）现有研究的局限与挑战（Gap 1）**
尽管关于DMA的理论研究已取得显著进展，但将理论转化为实际工程应用仍面临严峻挑战，主要体现在以下两个方面：

首先，**理想连续相位假设与实际离散硬件约束存在鸿沟**。在工程实现中，为了降低控制电路的复杂度和成本，DMA单元通常采用1-bit（开/关）或低比特（如2-bit）的数字控制方案，即单元的相位响应只能取有限个离散值。现有的基于连续相位的算法直接应用于离散系统时，需要进行强制量化（Rounding），这将破坏波束赋形的相干性，导致严重的天线增益损失和旁瓣抬升。

其次，**缺乏针对离散与耦合约束的高效算法**。由于DMA单元遵循洛伦兹谐振模型，其**幅度和相位存在固有的非线性耦合**，且离散化后的波束赋形问题本质上是一个非凸的整数规划问题（Integer Programming）。现有的穷举搜索法随单元数量呈指数级复杂度，难以应用于大规模阵列；而传统的松弛算法在处理强耦合的离散约束时往往陷入局部最优。因此，**如何在考虑洛伦兹耦合约束和离散相位约束的前提下，设计低复杂度的波束赋形算法，并从理论上量化分析离散化带来的性能损失，是DMA技术走向实用化亟待解决的关键问题**。

---

 **参考文献标注建议：**

* [1] Sleasman T, Imani M F, Gollub J N, et al. Dynamic metasurface antennas as a new paradigm for radio detection and ranging[J]. Applied Physics Letters, 2015. (Smith团队的基础性论文)
* [2] Pulido-Mancera L, Sleasman T, Smith D R, et al. Analytical modeling of dynamic metasurface antennas[J]. IEEE Transactions on Antennas and Propagation, 2021. (洛伦兹模型的详细分析)
* [3] [插入你查阅的关于DMA连续相位优化的相关文献]
* [4] [插入你查阅的关于全息波束赋形算法的相关文献]

 **写作提示：**

这段文字明确了Gap 1：**“连续假设 vs. 离散现实”** 以及 **“幅相独立 vs. 洛伦兹耦合”**。这不仅指出了前人的不足，更直接为你第二章（性能极限分析）和第三章（离散算法设计）提供了最为坚实的立论依据。

\subsection{捏合天线系统研究现状}

介绍第一次提出NTT DOCOMO

介绍性能分析，介绍波束赋形，介绍定位，介绍通感或者感知。

介绍挑战






\begin{figure}[!htbp]
	\centering
	\includegraphics[width=0.6\linewidth]{intro/intro_PASS.pdf}
	\caption{使用捏合天线部署通信区域\cite{fukuda2022pinching}}
	\label{fig:intro_PASS}
\end{figure}


\begin{figure}[!htbp]
	\centering
	\includegraphics[width=0.7\linewidth]{intro/intro_PASS_demo.pdf}
	\caption{通过捏合天线进行视频传输实验\cite{fukuda2022pinching}}
	\label{fig:intro_PASS_demo}
\end{figure}

\subsection{可重构智能表面研究现状}

介绍基本概念第一次提出、崔铁军和RFocus

介绍波束赋形算法

介绍系统研制

介绍不足

\subsection{现有研究中存在的不足}

DMA的理论不完善

从体系结构角度看，DMA通过在发射端引入电磁响应可重构的超材料单元，显著降低了射频链路数量，为大规模阵列的低功耗实现提供了新的硬件基础。然而，这种结构性简化也使得波束赋形过程受到波导传播效应与单元耦合特性的联合约束，其可实现的波束空间及性能边界有别于传统相控阵，亟需从理论层面对其波束赋形机理与性能极限进行系统分析。

因此，可重构天线系统的性能分析与波束赋形算法设计仍然有许多尚未解决的问题

\section{研究内容与章节安排}\label{sec:content_and_arrangement}

\subsection{论文主要研究内容}\label{sec:content}

本文主要研究xxx

\subsection{论文章节安排}\label{sec:arrangement}

全文共分为7章。主要内容如下：

第1章 绪论：简要介绍论文的研究背景、国内外研究现状、存在的问题，给出全文的主要研究内容。为了让读者更容易理解全文，建议用一个文档结构图给出各章节逻辑关系。

第2章 样式使用说明：一般应介绍该章的主要内容，说明这一章做什么。

第3章 实验研究类论文：……

第4章 理论或算法类论文的结构：……

第5章 总结与展望：给出全文的主要内容及结论，总结本文的创新点，并对未来的工作进行展望。





